% SKY2100 -- Exercise 9 (self-contained article; no external .sty needed)
\documentclass[11pt]{article}

% ---- Packages ------------------------------------------------------
\usepackage[a4paper,margin=2.2cm]{geometry}
\usepackage[T1]{fontenc}
\usepackage[utf8]{inputenc}
\usepackage{xcolor}
\usepackage{enumitem}
\usepackage{pgffor}
\usepackage{amssymb}
\usepackage{array}

% ---- Colours -------------------------------------------------------
\definecolor{kred}{HTML}{D81E2E}
\definecolor{kgrey}{HTML}{595959}

\usepackage[colorlinks=true,urlcolor=kred,linkcolor=kred]{hyperref}

% ---- Worksheet macros (inlined) ------------------------------------
\newcommand{\wbsection}[1]{\par\vspace{11pt}\noindent
  {\large\bfseries\color{kred}#1}\par\vspace{2pt}
  {\color{kred!35}\hrule height 1pt}\par\vspace{6pt}}

\newcommand{\task}[1]{\par\vspace{4pt}\noindent
  \textcolor{kred}{\textbf{$\triangleright$}}~#1\par\vspace{2pt}}

% Answers are discussed verbally -- no writing rules, just a small gap.
\newcommand{\answerlines}[1]{\par\vspace{0.4em}}

\newcommand{\boxlabel}[2]{\par\vspace{3pt}\noindent
  \fbox{\parbox{\dimexpr\linewidth-2\fboxsep-2\fboxrule}{\textbf{#1:}\enspace #2}}\par\vspace{3pt}}
\newcommand{\evidence}[1]{\boxlabel{Evidence to capture}{#1}}
\newcommand{\note}[1]{\boxlabel{Note}{#1}}
\newcommand{\caution}[1]{\boxlabel{Caution}{#1}}

% ---- Aha box + no italics (added) ----------------------------------
\newcommand{\aha}[1]{\par\vspace{5pt}\noindent
  \setlength{\fboxrule}{1.2pt}%
  {\color{kred}\fbox{\normalcolor\parbox{\dimexpr\linewidth-2\fboxsep-2\fboxrule}{%
     {\bfseries\color{kred}AHA.}\enspace #1}}}%
  \setlength{\fboxrule}{0.4pt}\par\vspace{5pt}}
\renewcommand{\emph}[1]{\textup{#1}}

\newcommand{\cmd}[1]{\texttt{#1}}
\newcommand{\cbox}{$\square$\enspace}

% ---- Hint and solutions divider -- same definitions as in kristiania-workbook.sty ----
\newcommand{\hint}[1]{\par\vspace{1pt}{\footnotesize\color{kgrey}\textbf{Hint:}\enspace #1}\par\vspace{2pt}}
\newcommand{\solheader}{\par\vspace{9pt}\noindent{\color{kred}\rule{\linewidth}{1.2pt}}\par\vspace{4pt}}

\newcommand{\signoff}{\par\vspace{9pt}\noindent
  {\color{kred}\rule{\linewidth}{0.4pt}}\par\vspace{3pt}
  \noindent{\footnotesize\color{kgrey}\textbf{Progress check:} Discuss your answers in the group, then
  talk the teacher through them verbally at the next session only to track progress; nothing is handed
  in, signed or graded.}\par}

\setlist[enumerate]{itemsep=3pt,topsep=2pt,parsep=0pt}
\setlist[itemize]{itemsep=2pt,topsep=2pt,parsep=0pt}
\setlength{\parindent}{0pt}
\setlength{\parskip}{2pt}

\begin{document}
\begin{center}
  {\LARGE\bfseries Watch your service --- find the attack in the logs}\\[3pt]
  {\large Session 9 -- Security monitoring \& IDS}\\[5pt]
  {\small Exercise 9 \textperiodcentered\ Session 9 lab \textperiodcentered\ Estimated time: 90--120 min}
\end{center}
\vspace{2pt}\hrule\vspace{1.1em}

Write your answers in the answer document \cmd{SKY2100\_Exercise9\_Answers.docx}. Your Nextcloud is encrypted (S6), hardened (S7) and shielded by a WAF (S8). Those
controls already \emph{produce} telemetry. Today you collect it in one place, then hunt for the attacks
you blocked last week and a brute-force pattern -- the difference between \emph{blocking} an attack and
\emph{seeing} it.

\noindent\textbf{Caution:} Monitor and generate traffic only on your \emph{own} VM. The ``attacks'' below
are harmless probes against your own service.

\wbsection{1\quad Collect the logs in one place}
\task{Identify and gather the security-relevant logs your VM already produces:}
\begin{quote}\ttfamily
sudo tail -n 20 /var/log/auth.log\\
sudo fail2ban-client status sshd\\
sudo tail -n 20 /var/log/modsec\_audit.log\hfill\# WAF (path may differ)\\
sudo journalctl -u nginx --no-pager | tail -n 20
\end{quote}
\note{No \cmd{/var/log/auth.log}? Ubuntu without rsyslog keeps these lines in the journal:
\cmd{sudo journalctl -u ssh --no-pager | tail -n 20}, and in Section~3 use
\cmd{sudo journalctl -u ssh | grep "Failed password"} instead of the \cmd{grep} on the file.}

\task{Which log would you check first if you suspected (a) a brute-force login, (b) a blocked web
attack, (c) a service crash? One line.}

\wbsection{2\quad Find last week's attack in the WAF log}
\task{Re-send one of your S8 test payloads, then locate it in the ModSecurity audit log:}
\begin{quote}\ttfamily
curl -k "https://<your-host>/?q=<script>alert(1)</script>"\\
sudo grep -iE "alert\textbackslash(1\textbackslash)|XSS|941" /var/log/modsec\_audit.log | tail
\end{quote}
\task{Which CRS rule fired (rule ID / message), what was the source IP, and what status was returned?}

\wbsection{3\quad Spot a pattern --- brute force over time}
\task{Generate a few failed SSH logins (from another host or with a wrong key/user), then look for the
\emph{pattern}, not the single line:}
\begin{quote}\ttfamily
sudo grep "Failed password" /var/log/auth.log | tail -n 20\\
sudo grep "Failed password" /var/log/auth.log | wc -l\\
sudo fail2ban-client status sshd
\end{quote}
\task{How many failures, from how many distinct IPs, over what time window? Why is the \emph{pattern} the
alert, not any single failed login?}

\aha{You blocked these attacks weeks ago, so why look at them now? Because blocking and seeing are
different things. fail2ban and the WAF stopped the attempts, but only the logs tell you they happened
at all, who tried, and how often. Notice too what counts as the signal. A single failed login is
nothing; people mistype passwords all the time. The alert is the pattern, many failures from one or a
few addresses in a short window, which no single line reveals. This is exactly why logs have to be
collected in one place and read together, and it is the groundwork for the timeline you build next week.}

\wbsection{4\quad (Optional) Run an IDS --- Suricata}
\task{Install Suricata, run it on your interface, and read one alert:}
\begin{quote}\ttfamily
sudo apt install -y suricata\\
sudo suricata -i <iface> \&\\
sudo tail -f /var/log/suricata/fast.log
\end{quote}
\task{Is Suricata here a NIDS or a HIDS, and is it signature- or anomaly-based by default? One line.}

\wbsection{5\quad Cloud transfer: Microsoft Learn}
\task{Work through ``Introduction to Microsoft Sentinel'' and note what a cloud-native SIEM does.}
\par\noindent\mbox{\href{https://learn.microsoft.com/training/modules/intro-to-azure-sentinel/}{learn.microsoft.com/training/modules/intro-to-azure-sentinel}}
\task{Name the three things a SIEM does with telemetry (hint: collect, \dots, \dots), and one advantage
of a SIEM over reading each host's logs by hand.}

\wbsection{6\quad Azure reflection}
\task{CSA distinguishes \emph{logs} (slow path) from \emph{events} (fast path). For detecting a
fast-moving management-plane attack, which would you rely on, and why?}
\task{Your fail2ban and WAF \emph{blocked} attacks. Why is also \emph{logging and alerting} on them
valuable, even though they were already stopped?}

\wbsection{7\quad Knowledge check}
\begin{enumerate}
\item Why is ``assume breach'' the starting point for monitoring?
\hint{Prevention is never perfect — plan to catch what gets through.}
\item State the difference between a \textbf{log} and an \textbf{event} (CSA).
\hint{One is the full, durable record; one is a single notable occurrence.}
\item State the difference between an \textbf{event} and an \textbf{alert}.
\hint{Raw captured data, versus an actionable notification derived from it.}
\item Name three telemetry sources you would collect on a VM.
\hint{auth.log, fail2ban, the WAF audit log, the web access log — any three.}
\item Why must logs be \emph{centralised} and \emph{protected}?
\hint{Scattered logs cannot be correlated; local ones an attacker can edit.}
\item Signature vs anomaly detection: one strength and one weakness of each.
\hint{One is precise but blind to new attacks; one catches new ones but is noisy.}
\item Difference between \textbf{NIDS} and \textbf{HIDS}, with one example tool each.
\hint{One watches network traffic; one watches a single host's logs and files.}
\item Difference between an \textbf{IDS} and an \textbf{IPS}.
\hint{One alerts, one blocks — a WAF is the blocking kind for HTTP.}
\item What is a \textbf{SIEM}, and what does ``correlation'' add?
\hint{It aggregates many sources; correlation ties separate events into one story.}
\item What does the \textbf{MITRE ATT\&CK} framework provide?
\hint{A map of real attacker tactics — context for which indicator means what.}
\item Define \textbf{MTTD} and why it matters.
\hint{Mean time to detect — the longer the dwell time, the worse the damage.}
\item What is \textbf{alert fatigue}, and how is it reduced?
\hint{Too many low-value alerts hide the real one; prioritise and tune.}
\end{enumerate}

\signoff

\clearpage
\solheader
{\LARGE\bfseries\color{kred} Solutions}\\[2pt]
{\small Work through the lab first, then check here. Model answers are indicative — equivalent reasoning is fine.}\par\vspace{8pt}
\noindent\textbf{Note:} Model answers are indicative -- accept equivalent reasoning. Multiple-choice
answers are in \textbf{bold}.

\wbsection{Knowledge check}
\begin{enumerate}
\item Prevention (firewall, TLS, hardening, WAF) is never perfect -- new vulnerabilities, stolen
      credentials and novel payloads slip through. ``Assume breach'' means you also build detection to
      catch what gets in, rather than trusting that nothing will.
\item A \textbf{log} is a relatively complete, detailed, durable record of activity (CRUD), often batched
      and slightly delayed. An \textbf{event} records only a change (C-UD), is ephemeral (not kept unless
      saved) and is available within seconds (CSA p.~125).
\item An \textbf{event} is raw captured data/activity; an \textbf{alert} is an \emph{actionable
      notification} derived from analysing events, signalling a potential threat. Events are the input;
      alerts are the output worth acting on (CSA p.~124).
\item Any three of: \cmd{/var/log/auth.log} (logins/sudo), fail2ban log (bans), the WAF/ModSecurity audit
      log (blocked attacks), the web access log (requests/status), \cmd{journalctl}/syslog (service \&
      system events). (CSA management-plane / service / resource logs, p.~128.)
\item \textbf{Centralised:} scattered logs cannot be correlated and give no overview of the whole estate
      (CSA p.~125). \textbf{Protected:} an attacker who can edit/delete logs erases their trail (the C-UD
      evasion, CSA p.~137); ship logs off-host, append-only.
\item \textbf{Signature:} precise, low false positives, but blind to novel attacks. \textbf{Anomaly:} can
      catch unknown attacks by learning ``normal'', but more false positives and needs tuning (CSA
      p.~126).
\item \textbf{NIDS} inspects network traffic (e.g.\ Suricata, Snort); \textbf{HIDS} watches one
      host's logs, files and processes (e.g.\ Wazuh, OSSEC). NIDS sees the knock; HIDS sees what happens
      inside.
\item An \textbf{IDS} detects and \emph{alerts}; an \textbf{IPS} detects and \emph{blocks}. The WAF is an
      application-layer IPS.
\item A \textbf{SIEM} aggregates, analyses and correlates logs/events from many sources (CSA p.~138).
      \textbf{Correlation} ties separate events (failed login + config change + data export) into one
      meaningful incident no single log would reveal.
\item \textbf{MITRE ATT\&CK} maps real attacker tactics and techniques, giving context for which
      indicators relate to which attack stage, so teams focus on the attacks most relevant to them (CSA
      p.~126).
\item \textbf{MTTD} = mean time to detect: the gap between an attack starting and being noticed. It
      matters because the longer the dwell time, the more an attacker learns, spreads and exfiltrates.
\item \textbf{Alert fatigue:} too many low-value alerts cause analysts to miss the real one. Reduced by
      prioritising (severity/confidence), tuning out false positives, and correlating many events into
      fewer incidents (CSA p.~124).
\end{enumerate}

\wbsection{Lab checkpoints}
\begin{itemize}
\item \textbf{Collect:} students should locate at least \cmd{auth.log}, the fail2ban status, the WAF
      audit log and the web/access log. First-check mapping: brute force $\rightarrow$
      \cmd{auth.log}/fail2ban; blocked web attack $\rightarrow$ WAF audit log; crash $\rightarrow$
      \cmd{journalctl}/syslog.
\item \textbf{WAF log:} the re-sent XSS probe should appear with a CRS rule (typically a 941xxx XSS rule),
      the source IP and a 403. This is the ``see what you blocked'' moment -- blocking and logging are
      different jobs.
\item \textbf{Brute-force pattern:} several \cmd{Failed password} lines from one or more IPs in a short
      window. The teaching point (CSA p.~137): one failure is noise; ``repeated login attempts from
      different IPs'' over time is the alert-worthy pattern.
\item \textbf{Suricata (optional):} by default a \textbf{NIDS}, \textbf{signature}-based (rule sets). A
      single \cmd{fast.log} alert (e.g.\ a scan or a rule hit) is enough.
\end{itemize}

\wbsection{Cloud transfer \& reflection}
\begin{itemize}
\item \textbf{SIEM:} collects, \emph{correlates} and \emph{alerts} on telemetry from many sources.
      Advantage over per-host logs: one place to search, cross-source correlation, and consistent
      alerting -- you see the whole estate, not one machine.
\item \textbf{Logs vs events for a fast attack:} rely on \textbf{events} (the fast path, near-real-time,
      e.g.\ Defender/GuardDuty/Sentinel) for management-plane attacks, because logs are the slow path and
      delays are unacceptable when attacks move fast (CSA p.~126, 137).
\item \textbf{Why log blocked attacks:} blocking stops \emph{this} attempt; logging/alerting reveals that
      you are being targeted, by whom and how often, feeds correlation and forensics, and meets
      compliance -- detection and prevention are complementary, not redundant.
\end{itemize}

\wbsection{References}
CSA Security Guidance v5, \textbf{Domain~6: Security Monitoring} (p.~124--143, verified): cloud monitoring
challenges and management plane (p.~125--126), logs vs events and their uses (p.~125--126), events vs
alerts (p.~124), timeliness and key indicators / MITRE ATT\&CK (p.~126), telemetry sources (p.~127--128),
detection paths and the C-UD evasion example (p.~137), SIEM vs CSP alerts vs CSPM (p.~138). Comer,
\emph{The Cloud Computing Book}, Ch.~17.8 -- AI \& security analytics: context-aware anomaly detection and
DDoS detection (p.~239), verified. MS Learn: \emph{Introduction to Microsoft Sentinel} (verified, June
2026) -- cloud-native SIEM + SOAR.

\end{document}
